
\documentclass{article}
\usepackage{colortbl}
\usepackage{makecell}
\usepackage{multirow}
\usepackage{supertabular}

\begin{document}

\newcounter{utterance}

\centering \large Interaction Transcript for game `dond', experiment `semi\_en', episode 5 with Qwen3{-}30B{-}A3B{-}t0.0.
\vspace{24pt}

{ \footnotesize  \setcounter{utterance}{1}
\setlength{\tabcolsep}{0pt}
\begin{supertabular}{c@{$\;$}|p{.15\linewidth}@{}p{.15\linewidth}p{.15\linewidth}p{.15\linewidth}p{.15\linewidth}p{.15\linewidth}}
    \# & $\;$A & \multicolumn{4}{c}{Game Master} & $\;\:$B\\
    \hline

    \theutterance \stepcounter{utterance}  
    & & \multicolumn{4}{p{0.6\linewidth}}{
        \cellcolor[rgb]{0.9,0.9,0.9}{
            \makecell[{{p{\linewidth}}}]{
                \texttt{\tiny{[P1$\langle$GM]}}
                \texttt{You are playing a negotiation game in which you have to agree on how to divide a set of items among you and another player.} \\
\\ 
\texttt{Rules:} \\
\texttt{(a) You and the other player are given a set of items. Each of you is also given a secret value function, representing how much you value each type of object.} \\
\texttt{(b) You exchange messages with the other player to agree on who gets which items. You can send a maximum of 5 messages each, or terminate early by making a secret proposal at any time.} \\
\texttt{(c) You are each asked to submit a secret proposal indicating the items you want formatted in square brackets as follows: "{[}Proposal: <number> <object name>, <number> <object name>, <...>{]}"} \\
\texttt{(d) If your proposals are complementary, i.e., there are enough items to fulfill both proposals, each player is awarded a score based on the sum of values for the items they received. Otherwise, both of you get zero points.} \\
\texttt{(e) Your goal is to maximize the score you receive.} \\
\\ 
\texttt{Let us start.} \\
\\ 
\texttt{The set of available items is:} \\
\\ 
\texttt{2 dice, 1 ball, 1 glue, 1 battery} \\
\\ 
\texttt{Your secret value function is:} \\
\\ 
\texttt{dice: 1, ball: 7, glue: 0, battery: 1} \\
\\ 
\texttt{IMPORTANT: Your messages, unless it is the secret proposal, are directly transmitted to the other player, so do not include any response to the rules or text announcing your message. To make a secret proposal, use the indicated format. Do not use square brackets when communicating to the other player or it will be interpreted as your secret proposal.} \\
            }
        }
    }
    & \\ \\

    \theutterance \stepcounter{utterance}  
    & \multicolumn{4}{p{0.6\linewidth}}{
        \cellcolor[rgb]{0.9,0.9,0.9}{
            \makecell[{{p{\linewidth}}}]{
                \texttt{\tiny{[P1$\rangle$GM]}}
                \texttt{<think>} \\
\texttt{Okay, let's see. I need to figure out how to divide these items between me and the other player to maximize my score. The items are 2 dice, 1 ball, 1 glue, 1 battery. My values are dice:1, ball:7, glue:0, battery:1. So the ball is really valuable to me, and the glue isn't worth anything. The other player's values aren't known, but I have to guess or find a way to communicate effectively.} \\
\\ 
\texttt{First, I should think about what the other player might want. Since I don't know their values, maybe I can infer based on the items. The ball is high value for me, so maybe they also value it? Or maybe they don't. But since I can't know, I need to make a proposal that's likely to be complementary.} \\
\\ 
\texttt{The key is that if our proposals are complementary, we both get points. Otherwise, zero. So I need to make a proposal that the other player can accept without conflicting. But how?} \\
\\ 
\texttt{Let me think about possible strategies. If I take the ball, that's 7 points for me. But if the other player also wants the ball, we can't both have it. So maybe I should try to get the ball, but also make sure that the other player gets items they value. But without knowing their values, it's tricky.} \\
\\ 
\texttt{Alternatively, maybe I can propose a division where I take the ball and the other player takes the dice and battery, which have lower value for me. But I need to check if that's possible. Wait, the other player's values are unknown. So maybe I should try to communicate my preferences in a way that the other player can respond.} \\
\\ 
\texttt{But since I can send up to 5 messages, maybe I can start by asking what they want. But the problem is that messages are directly transmitted. So if I ask, "What do you want?" the other player would get that message. But maybe they would respond. However, I need to be careful not to reveal too much.} \\
\\ 
\texttt{Alternatively, maybe I can make a proposal that includes the ball and some other items, but also leaves room for the other player to get their items. For example, if I propose to take the ball and the battery, and the other player can take the dice and glue. But the glue is worth 0 to me, so maybe that's okay. But the other player might} \\
            }
        }
    }
    & & \\ \\

    \theutterance \stepcounter{utterance}  
    & & & \multicolumn{2}{p{0.3\linewidth}}{
        \cellcolor[rgb]{0.9,0.9,0.9}{
            \makecell[{{p{\linewidth}}}]{
                \texttt{\tiny{[GM$|$GM]}}
                \texttt{continue} \\
            }
        }
    }
    & & \\ \\

    \theutterance \stepcounter{utterance}  
    & & & \multicolumn{2}{p{0.3\linewidth}}{
        \cellcolor[rgb]{0.9,0.9,0.9}{
            \makecell[{{p{\linewidth}}}]{
                \texttt{\tiny{[GM$|$GM]}}
                \texttt{<think>} \\
\texttt{Okay, let's see. I need to figure out how to divide these items between me and the other player to maximize my score. The items are 2 dice, 1 ball, 1 glue, 1 battery. My values are dice:1, ball:7, glue:0, battery:1. So the ball is really valuable to me, and the glue isn't worth anything. The other player's values aren't known, but I have to guess or find a way to communicate effectively.} \\
\\ 
\texttt{First, I should think about what the other player might want. Since I don't know their values, maybe I can infer based on the items. The ball is high value for me, so maybe they also value it? Or maybe they don't. But since I can't know, I need to make a proposal that's likely to be complementary.} \\
\\ 
\texttt{The key is that if our proposals are complementary, we both get points. Otherwise, zero. So I need to make a proposal that the other player can accept without conflicting. But how?} \\
\\ 
\texttt{Let me think about possible strategies. If I take the ball, that's 7 points for me. But if the other player also wants the ball, we can't both have it. So maybe I should try to get the ball, but also make sure that the other player gets items they value. But without knowing their values, it's tricky.} \\
\\ 
\texttt{Alternatively, maybe I can propose a division where I take the ball and the other player takes the dice and battery, which have lower value for me. But I need to check if that's possible. Wait, the other player's values are unknown. So maybe I should try to communicate my preferences in a way that the other player can respond.} \\
\\ 
\texttt{But since I can send up to 5 messages, maybe I can start by asking what they want. But the problem is that messages are directly transmitted. So if I ask, "What do you want?" the other player would get that message. But maybe they would respond. However, I need to be careful not to reveal too much.} \\
\\ 
\texttt{Alternatively, maybe I can make a proposal that includes the ball and some other items, but also leaves room for the other player to get their items. For example, if I propose to take the ball and the battery, and the other player can take the dice and glue. But the glue is worth 0 to me, so maybe that's okay. But the other player might} \\
            }
        }
    }
    & & \\ \\

    \theutterance \stepcounter{utterance}  
    & & \multicolumn{4}{p{0.6\linewidth}}{
        \cellcolor[rgb]{0.9,0.9,0.9}{
            \makecell[{{p{\linewidth}}}]{
                \texttt{\tiny{[GM$\rangle$P2]}}
                \texttt{You are playing a negotiation game in which you have to agree on how to divide a set of items among you and another player.} \\
\\ 
\texttt{Rules:} \\
\texttt{(a) You and the other player are given a set of items. Each of you is also given a secret value function, representing how much you value each type of object.} \\
\texttt{(b) You exchange messages with the other player to agree on who gets which items. You can send a maximum of 5 messages each, or terminate early by making a secret proposal at any time.} \\
\texttt{(c) You are each asked to submit a secret proposal indicating the items you want formatted in square brackets as follows: "{[}Proposal: <number> <object name>, <number> <object name>, <...>{]}"} \\
\texttt{(d) If your proposals are complementary, i.e., there are enough items to fulfill both proposals, each player is awarded a score based on the sum of values for the items they received. Otherwise, both of you get zero points.} \\
\texttt{(e) Your goal is to maximize the score you receive.} \\
\\ 
\texttt{Let us start.} \\
\\ 
\texttt{The set of available items is:} \\
\\ 
\texttt{2 dice, 1 ball, 1 glue, 1 battery} \\
\\ 
\texttt{Your secret value function is:} \\
\\ 
\texttt{dice: 0, ball: 2, glue: 8, battery: 0} \\
\\ 
\texttt{IMPORTANT: Your messages, unless it is the secret proposal, are directly transmitted to the other player, so do not include any response to the rules or text announcing your message. To make a secret proposal, use the indicated format. Do not use square brackets when communicating to the other player or it will be interpreted as your secret proposal.} \\
\\ 
\\ 
\texttt{<think>} \\
\texttt{Okay, let's see. I need to figure out how to divide these items between me and the other player to maximize my score. The items are 2 dice, 1 ball, 1 glue, 1 battery. My values are dice:1, ball:7, glue:0, battery:1. So the ball is really valuable to me, and the glue isn't worth anything. The other player's values aren't known, but I have to guess or find a way to communicate effectively.} \\
\\ 
\texttt{First, I should think about what the other player might want. Since I don't know their values, maybe I can infer based on the items. The ball is high value for me, so maybe they also value it? Or maybe they don't. But since I can't know, I need to make a proposal that's likely to be complementary.} \\
\\ 
\texttt{The key is that if our proposals are complementary, we both get points. Otherwise, zero. So I need to make a proposal that the other player can accept without conflicting. But how?} \\
\\ 
\texttt{Let me think about possible strategies. If I take the ball, that's 7 points for me. But if the other player also wants the ball, we can't both have it. So maybe I should try to get the ball, but also make sure that the other player gets items they value. But without knowing their values, it's tricky.} \\
\\ 
\texttt{Alternatively, maybe I can propose a division where I take the ball and the other player takes the dice and battery, which have lower value for me. But I need to check if that's possible. Wait, the other player's values are unknown. So maybe I should try to communicate my preferences in a way that the other player can respond.} \\
\\ 
\texttt{But since I can send up to 5 messages, maybe I can start by asking what they want. But the problem is that messages are directly transmitted. So if I ask, "What do you want?" the other player would get that message. But maybe they would respond. However, I need to be careful not to reveal too much.} \\
\\ 
\texttt{Alternatively, maybe I can make a proposal that includes the ball and some other items, but also leaves room for the other player to get their items. For example, if I propose to take the ball and the battery, and the other player can take the dice and glue. But the glue is worth 0 to me, so maybe that's okay. But the other player might} \\
            }
        }
    }
    & \\ \\

    \theutterance \stepcounter{utterance}  
    & & & \multicolumn{4}{p{0.6\linewidth}}{
        \cellcolor[rgb]{0.9,0.9,0.9}{
            \makecell[{{p{\linewidth}}}]{
                \texttt{\tiny{[GM$\langle$P2]}}
                \texttt{<think>} \\
\texttt{Okay, let's try to think through this step by step. The goal is to maximize my score, which depends on the items I get based on my value function. My values are: dice (0), ball (2), glue (8), battery (0). Wait, no, the user said "Your secret value function is: dice: 0, ball: 2, glue: 8, battery: 0." Wait, that's conflicting with the initial message. Wait, the user wrote:} \\
\\ 
\texttt{"Your secret value function is:} \\
\\ 
\texttt{dice: 0, ball: 2, glue: 8, battery: 0"} \\
\\ 
\texttt{Wait, but in the initial problem statement, the user said:} \\
\\ 
\texttt{"Your secret value function is:} \\
\\ 
\texttt{dice: 0, ball: 2, glue: 8, battery: 0"} \\
\\ 
\texttt{Wait, but in the problem statement, the user might have made a typo. Wait, the original problem says:} \\
\\ 
\texttt{"Your secret value function is:} \\
\\ 
\texttt{dice: 0, ball: 2, glue: 8, battery: 0"} \\
\\ 
\texttt{But then in the initial message, the user wrote:} \\
\\ 
\texttt{"Let us start.} \\
\\ 
\texttt{The set of available items is:} \\
\\ 
\texttt{2 dice, 1 ball, 1 glue, 1 battery} \\
\\ 
\texttt{Your secret value function is:} \\
\\ 
\texttt{dice: 0, ball: 2, glue: 8, battery: 0"} \\
\\ 
\texttt{Wait, so the user's value for ball is 2, glue is 8, dice and battery are 0. So the ball is worth 2, glue is 8. So the most valuable item for me is glue (8), then ball (2), and dice and battery are 0. So I should prioritize getting the glue and the ball. But the other player's values are unknown.} \\
\\ 
\texttt{But the problem is that if I propose a division, the other player's proposal must be complementary. So if I take the glue and the ball, the other player must get the remaining items (2 dice and 1 battery). But if the other player's proposal includes items that don't conflict, then we both get points.} \\
\\ 
\texttt{But how to ensure that? Since I don't know their values, maybe I should try to get the items that are most valuable to me, and see if the other player can get the rest.} \\
\\ 
\texttt{But the other player might also want the glue or the ball. So maybe I need to communicate in a way that the} \\
            }
        }
    }
     \\ \\

    \theutterance \stepcounter{utterance}  
    & & & \multicolumn{2}{p{0.3\linewidth}}{
        \cellcolor[rgb]{0.9,0.9,0.9}{
            \makecell[{{p{\linewidth}}}]{
                \texttt{\tiny{[GM$|$GM]}}
                \texttt{continue} \\
            }
        }
    }
    & & \\ \\

    \theutterance \stepcounter{utterance}  
    & & & \multicolumn{2}{p{0.3\linewidth}}{
        \cellcolor[rgb]{0.9,0.9,0.9}{
            \makecell[{{p{\linewidth}}}]{
                \texttt{\tiny{[GM$|$GM]}}
                \texttt{<think>} \\
\texttt{Okay, let's try to think through this step by step. The goal is to maximize my score, which depends on the items I get based on my value function. My values are: dice (0), ball (2), glue (8), battery (0). Wait, no, the user said "Your secret value function is: dice: 0, ball: 2, glue: 8, battery: 0." Wait, that's conflicting with the initial message. Wait, the user wrote:} \\
\\ 
\texttt{"Your secret value function is:} \\
\\ 
\texttt{dice: 0, ball: 2, glue: 8, battery: 0"} \\
\\ 
\texttt{Wait, but in the initial problem statement, the user said:} \\
\\ 
\texttt{"Your secret value function is:} \\
\\ 
\texttt{dice: 0, ball: 2, glue: 8, battery: 0"} \\
\\ 
\texttt{Wait, but in the problem statement, the user might have made a typo. Wait, the original problem says:} \\
\\ 
\texttt{"Your secret value function is:} \\
\\ 
\texttt{dice: 0, ball: 2, glue: 8, battery: 0"} \\
\\ 
\texttt{But then in the initial message, the user wrote:} \\
\\ 
\texttt{"Let us start.} \\
\\ 
\texttt{The set of available items is:} \\
\\ 
\texttt{2 dice, 1 ball, 1 glue, 1 battery} \\
\\ 
\texttt{Your secret value function is:} \\
\\ 
\texttt{dice: 0, ball: 2, glue: 8, battery: 0"} \\
\\ 
\texttt{Wait, so the user's value for ball is 2, glue is 8, dice and battery are 0. So the ball is worth 2, glue is 8. So the most valuable item for me is glue (8), then ball (2), and dice and battery are 0. So I should prioritize getting the glue and the ball. But the other player's values are unknown.} \\
\\ 
\texttt{But the problem is that if I propose a division, the other player's proposal must be complementary. So if I take the glue and the ball, the other player must get the remaining items (2 dice and 1 battery). But if the other player's proposal includes items that don't conflict, then we both get points.} \\
\\ 
\texttt{But how to ensure that? Since I don't know their values, maybe I should try to get the items that are most valuable to me, and see if the other player can get the rest.} \\
\\ 
\texttt{But the other player might also want the glue or the ball. So maybe I need to communicate in a way that the} \\
            }
        }
    }
    & & \\ \\

    \theutterance \stepcounter{utterance}  
    & & \multicolumn{4}{p{0.6\linewidth}}{
        \cellcolor[rgb]{0.9,0.9,0.9}{
            \makecell[{{p{\linewidth}}}]{
                \texttt{\tiny{[P1$\langle$GM]}}
                \texttt{<think>} \\
\texttt{Okay, let's try to think through this step by step. The goal is to maximize my score, which depends on the items I get based on my value function. My values are: dice (0), ball (2), glue (8), battery (0). Wait, no, the user said "Your secret value function is: dice: 0, ball: 2, glue: 8, battery: 0." Wait, that's conflicting with the initial message. Wait, the user wrote:} \\
\\ 
\texttt{"Your secret value function is:} \\
\\ 
\texttt{dice: 0, ball: 2, glue: 8, battery: 0"} \\
\\ 
\texttt{Wait, but in the initial problem statement, the user said:} \\
\\ 
\texttt{"Your secret value function is:} \\
\\ 
\texttt{dice: 0, ball: 2, glue: 8, battery: 0"} \\
\\ 
\texttt{Wait, but in the problem statement, the user might have made a typo. Wait, the original problem says:} \\
\\ 
\texttt{"Your secret value function is:} \\
\\ 
\texttt{dice: 0, ball: 2, glue: 8, battery: 0"} \\
\\ 
\texttt{But then in the initial message, the user wrote:} \\
\\ 
\texttt{"Let us start.} \\
\\ 
\texttt{The set of available items is:} \\
\\ 
\texttt{2 dice, 1 ball, 1 glue, 1 battery} \\
\\ 
\texttt{Your secret value function is:} \\
\\ 
\texttt{dice: 0, ball: 2, glue: 8, battery: 0"} \\
\\ 
\texttt{Wait, so the user's value for ball is 2, glue is 8, dice and battery are 0. So the ball is worth 2, glue is 8. So the most valuable item for me is glue (8), then ball (2), and dice and battery are 0. So I should prioritize getting the glue and the ball. But the other player's values are unknown.} \\
\\ 
\texttt{But the problem is that if I propose a division, the other player's proposal must be complementary. So if I take the glue and the ball, the other player must get the remaining items (2 dice and 1 battery). But if the other player's proposal includes items that don't conflict, then we both get points.} \\
\\ 
\texttt{But how to ensure that? Since I don't know their values, maybe I should try to get the items that are most valuable to me, and see if the other player can get the rest.} \\
\\ 
\texttt{But the other player might also want the glue or the ball. So maybe I need to communicate in a way that the} \\
            }
        }
    }
    & \\ \\

    \theutterance \stepcounter{utterance}  
    & \multicolumn{4}{p{0.6\linewidth}}{
        \cellcolor[rgb]{0.9,0.9,0.9}{
            \makecell[{{p{\linewidth}}}]{
                \texttt{\tiny{[P1$\rangle$GM]}}
                \texttt{<think>} \\
\texttt{Okay, let's start over. The user's value function is: dice: 0, ball: 2, glue: 8, battery: 0. So the most valuable items for me are glue (8) and ball (2). The other player's values are unknown, but I need to negotiate a division where both proposals are complementary.} \\
\\ 
\texttt{First, I need to figure out what the other player might want. Since I don't know their values, I have to make assumptions. Maybe they also value the glue or the ball. But without knowing, I need to make a proposal that's likely to be compatible.} \\
\\ 
\texttt{If I take the glue and the ball, that's 8 + 2 = 10 points for me. The remaining items are 2 dice and 1 battery. If the other player can take those, then our proposals are complementary. But if they also want the glue or ball, there's a conflict.} \\
\\ 
\texttt{Alternatively, maybe I should propose to take the glue and one die, leaving the ball and battery for the other player. But since the ball is worth 2 to me, maybe I should prioritize that. But I don't know.} \\
\\ 
\texttt{Another approach: since the other player's values are unknown, I can try to communicate my preferences in a way that they can respond. For example, if I send a message saying I want the glue and the ball, they might respond with their own preferences. But since messages are direct, I need to be careful not to reveal too much.} \\
\\ 
\texttt{But the problem is that I can only send 5 messages. So maybe the best strategy is to make a secret proposal that includes the most valuable items for me, hoping that the other player's proposal doesn't conflict.} \\
\\ 
\texttt{So my secret proposal would be to take the glue and the ball. That's {[}Proposal: 1 glue, 1 ball{]}. Then, the other player would have to take the remaining items: 2 dice and 1 battery. If their proposal doesn't conflict, we both get points. But if they also want the glue or ball, there's a conflict.} \\
\\ 
\texttt{Alternatively, maybe I should split the items differently. For example, take the glue and one die, and leave the ball and battery for the other player. But since the ball is worth 2 to me, maybe I should prioritize that.} \\
\\ 
\texttt{Wait, but the other player might not want the ball. If I take the ball, they might not have a} \\
            }
        }
    }
    & & \\ \\

    \theutterance \stepcounter{utterance}  
    & & & \multicolumn{2}{p{0.3\linewidth}}{
        \cellcolor[rgb]{0.9,0.9,0.9}{
            \makecell[{{p{\linewidth}}}]{
                \texttt{\tiny{[GM$|$GM]}}
                \texttt{proposal} \\
            }
        }
    }
    & & \\ \\

    \theutterance \stepcounter{utterance}  
    & & & \multicolumn{2}{p{0.3\linewidth}}{
        \cellcolor[rgb]{0.9,0.9,0.9}{
            \makecell[{{p{\linewidth}}}]{
                \texttt{\tiny{[GM$|$GM]}}
                [0, 1, 1, 0]
            }
        }
    }
    & & \\ \\

    \theutterance \stepcounter{utterance}  
    & & \multicolumn{4}{p{0.6\linewidth}}{
        \cellcolor[rgb]{0.9,0.9,0.9}{
            \makecell[{{p{\linewidth}}}]{
                \texttt{\tiny{[GM$\rangle$P2]}}
                \texttt{The other player has submitted a secret proposal. It is now your turn to submit a proposal of your own.} \\
            }
        }
    }
    & \\ \\

    \theutterance \stepcounter{utterance}  
    & & & \multicolumn{4}{p{0.6\linewidth}}{
        \cellcolor[rgb]{0.9,0.9,0.9}{
            \makecell[{{p{\linewidth}}}]{
                \texttt{\tiny{[GM$\langle$P2]}}
                \texttt{<think>} \\
\texttt{Okay, the other player has already submitted their secret proposal, and now it's my turn. I need to figure out what my proposal should be. But wait, the user hasn't provided the other player's proposal yet. Hmm, maybe I need to check the history.} \\
\\ 
\texttt{Wait, looking back at the conversation, the user said: "The other player has submitted a secret proposal. It is now your turn to submit a proposal of your own." But there's no previous message from the other player. That's confusing. Maybe there was a mistake in the setup. Alternatively, perhaps the user is expecting me to generate a proposal based on the given information, assuming that the other player's proposal is something that I need to respond to.} \\
\\ 
\texttt{But without knowing the other player's proposal, how can I determine if our proposals are complementary? This seems like a problem. Maybe the user intended for me to proceed with my own proposal, assuming that the other player's proposal is something that can be complementary. But that's unclear.} \\
\\ 
\texttt{Alternatively, maybe the user is testing my ability to handle incomplete information. In that case, I need to make a proposal that maximizes my chances of complementary proposals. Given that, I should consider what items I want and what the other player might want.} \\
\\ 
\texttt{My value function is: dice (0), ball (2), glue (8), battery (0). So I value glue the most, then ball. The other player might value different items. If I propose to take the glue and the ball, that would give me 10 points. But if the other player also wants the glue or the ball, there could be a conflict. However, if they don't, then it's complementary.} \\
\\ 
\texttt{Alternatively, maybe I should propose to take the glue and one die, leaving the ball and battery for the other player. But since dice and battery are worth 0 to me, it's better to take the glue and ball. But I need to ensure that the other player can get their items.} \\
\\ 
\texttt{But without knowing their proposal, it's risky. However, since I have to submit a proposal, I'll proceed with the one that gives me the highest possible score, assuming that the other player's proposal doesn't conflict. So my proposal would be to take the glue and the ball, which are the highest value items for me. That would be {[}Proposal: 1 glue, 1 ball{]}. But wait, there are 2 dice, 1} \\
            }
        }
    }
     \\ \\

    \theutterance \stepcounter{utterance}  
    & & & \multicolumn{2}{p{0.3\linewidth}}{
        \cellcolor[rgb]{0.9,0.9,0.9}{
            \makecell[{{p{\linewidth}}}]{
                \texttt{\tiny{[GM$|$GM]}}
                \texttt{proposal} \\
            }
        }
    }
    & & \\ \\

    \theutterance \stepcounter{utterance}  
    & & & \multicolumn{2}{p{0.3\linewidth}}{
        \cellcolor[rgb]{0.9,0.9,0.9}{
            \makecell[{{p{\linewidth}}}]{
                \texttt{\tiny{[GM$|$GM]}}
                [0, 1, 1, 0]
            }
        }
    }
    & & \\ \\

    \theutterance \stepcounter{utterance}  
    & & & \multicolumn{2}{p{0.3\linewidth}}{
        \cellcolor[rgb]{0.9,0.9,0.9}{
            \makecell[{{p{\linewidth}}}]{
                \texttt{\tiny{[GM$|$GM]}}
                [[0, 1, 1, 0], [0, 1, 1, 0]]
            }
        }
    }
    & & \\ \\

\end{supertabular}
}

\end{document}
